\documentclass{article}
\usepackage{hyperref}
\title{A normal status need not mean the hardware was sampled}
\begin{document}
\maketitle

The pinned \href{https://github.com/ros-controls/ros2_control/blob/0c94b008a82731967ba89e2e7379ec6d2f07580b/hardware_interface/src/resource_manager.cpp}{ros2\_control resource manager}
can return the same result and empty failure list after executing
both devices' reads or after executing neither. A state change for one
device acquires manager-wide locks. During its blocked lifecycle
operation, another device's read error need not be observed even if
its own device mutex is free and it shares no hardware group with
the device being reconfigured.

\paragraph{Mechanism.}
The lifted input to a recovery decision is not just the enum and
failure list. The pertinent fact is whether each required device
was sampled in this cycle and, where relevant, whether the selected
output and limit path ran. The native read and write initialize the
reported status to \texttt{OK} and clear the failed-device names
before trying the manager resource mutex. On failure to acquire the
mutex, each returns those initialized fields without invoking any
device. Joint-limit enforcement independently tries a different
mutex, but the resource manager's state-change operation holds both
of them around the lifecycle transition. At an individual device
mutex, the read loop skips just that device and can continue to
another. A true report needs the execution side of the relevant
path, not only the absence of a reported error.

\paragraph{Replays.}
The upstream source pin and blob are in \texttt{SOURCE.tsv}. The
harness executes the four original functions and native group-state
rule with two separate stub devices. Their callbacks have explicit
read, write and limiter counters. A condition variable pauses the
stub device A lifecycle transition after the original
\texttt{set\_component\_state} acquired the two locks. Device B
would return an error if read. Three cycles return \texttt{OK} with
empty lists for both read and write, execute neither device's
callbacks, and do not enforce any modeled limits. The first read
after releasing the old transition reports B's error and removes
its interfaces from availability in that replay. A matched
healthy cycle has the same read/write status pair but genuinely
calls both devices. The limit result is distinct in this fixture;
the downstream monitor considered here sees the read or write
result, not a full trace of all three operations.

In the scope control, holding only A's device mutex leaves B's
callbacks available; B's read error is reported while A is
skipped. This makes the shared manager mutex, rather than
a shared device or a shared hardware group, the relevant
dependency for the earlier trace. The code-selected device
operations and their order were checked. The lifecycle duration,
read result, limiter effects, and device internals are fixture
inputs. No hardware, simulation at real time scale, or installed
driver was exercised.

\paragraph{Consequence for the switch comparison.}
The local read error in the
\href{../local-fallback/Local_Fallback_Before_Switch_Return.tex}{fallback
study} cannot be credited during the masked cycle merely because the
failing device has a nominally independent scope. A consumer that uses
these read/write fields alone cannot distinguish the pair above;
a trusted completion record for the relevant device and cycle
separates them. This improves the correctness of the evidence,
not the timeliness of the underlying observation. Any bounded
reaction claim also requires a bound on the lock and subsequent
sampling delays, or a different authority and observation path
that stays available when this transition stalls. The replay makes
no claim about whether a real device holds, drops, or limits
the last command through the skipped cycles. This is a stock
baseline limitation for I1, not an architecture-level pass
or an assessment of a deployed fleet.
\end{document}
